% Revisão consolidada: 25/11/2026. Não antecipa resultados finais.
\chapter{Fundamentação teórica}
\section{Polifarmácia, adesão e rastreabilidade}
O número de medicamentos é uma forma frequente de caracterizar polifarmácia, mas não descreve sozinho a adequação da prescrição, a complexidade do regime nem as necessidades individuais. A revisão de Masnoon e colaboradores identifica variação importante nas definições. Neste trabalho, polifarmácia contextualiza o público pretendido, sem constituir critério de inclusão de participantes ainda não definidos. \cite{MASNOON2017}
Adesão, disponibilidade de uma dose, retirada de um recipiente e ingestão são eventos diferentes. O registro de uma etapa não autoriza assumir as demais. A avaliação deve definir precisamente qual evento foi observado, por qual instrumento, em que janela temporal e com qual incerteza. Essa distinção orienta a separação de eventos do dispositivo e inferências visuais na documentação do DoseCerta. \cite{DC_COMPILADO}
\section{Dispensação unitária e mecanismos rotativos}
Um mecanismo de disco com cavidades busca capturar unidades, transportá-las e liberá-las por uma abertura de saída. A geometria precisa conciliar dimensões do comprimido, tolerância de fabricação, retenção durante o percurso e liberação na posição desejada. O modelo recuperado do DoseCerta descreve disco com cinco cavidades, reservatório e placa de apoio. A correspondência entre geometria digital e peça fabricada exige medição. \cite{DC_COMPILADO}
A inclinação adotada nesta revisão é 35\ensuremath{^{\circ}}, conforme decisão da autora e pré-relatório. Ela é um parâmetro de projeto, não um resultado de ensaio. A documentação identifica outra inclinação no modelo V75; por isso, o CAD recuperado não deve ser apresentado automaticamente como geometria final. \cite{DC_PREREL,DC_COMPILADO}
\section{Motor de passo e estágio de potência}
A sequência de excitação das bobinas produz deslocamentos discretos. A redução mecânica e o regime de acionamento afetam velocidade, torque, folga e repetibilidade. Contar fases comandadas não demonstra que o eixo ou o comprimido percorreu o movimento esperado. No projeto, o conjunto documentado é 28BYJ-48 de 5 V com placa ULN2003, alimentado por ramal próprio e não por GPIO. \cite{DC_CIRCUITOS}
A conversão entre fases e ângulo requer referência do conjunto instalado e medição de uma volta, sem assumir uma razão nominal como calibração. Limites locais de duração e de número de fases reduzem a dependência do servidor para interromper uma operação, mas não substituem a verificação do mecanismo. \cite{DC_COMPILADO}
\section{Sensor óptico reflexivo}
O TCRT5000 integra emissor infravermelho e fototransistor, com resposta dependente da distância e da refletância do alvo. A folha de dados o caracteriza como sensor reflexivo, não como instrumento de identificação química ou de confirmação clínica. Cor, material, posição e luz ambiente interferem na interpretação de leituras. \cite{VISHAY_TCRT}
Na topologia descrita para o DoseCerta, o coletor é elevado a 3,3 V por resistor e o emissor é ligado ao terra; maior reflexão tende a reduzir a tensão no ADC. Essa leitura serve à observação da cavidade. Não se pode rebatizá-la como barreira óptica na saída, nem usar uma cavidade vazia como prova isolada de ingestão. \cite{DC_CIRCUITOS,DC_COMPILADO}
\section{Arquitetura distribuída e comunicação I²C}
A comunicação I²C utiliza sinais de dados e relógio compartilhados, com saídas de dreno/coletor aberto e elevação de nível apropriada. Endereçamento distinto, capacitância, resistores equivalentes e tempos de subida condicionam o funcionamento elétrico; o tamanho físico da montagem e seus contatos fazem parte dessa análise. \cite{NXP_I2C}
O CRC detecta determinadas corrupções acidentais, mas não autentica a origem de uma mensagem. Identificadores de comando, estados de inicialização e resultados persistidos atendem a outra necessidade: reconhecer reenvios e impedir que uma perda de resposta provoque um segundo movimento. A documentação do protocolo DC1 aplica esses mecanismos em quadros de tamanho fixo. \cite{DC_COMPILADO}
\section{Persistência, idempotência e operação local}
Idempotência significa que a repetição de uma solicitação não deve produzir novo efeito indevido. No controle físico, uma resposta perdida pode deixar o resultado desconhecido; repetir o movimento não é uma estratégia segura de confirmação. O projeto prevê consulta de estado e reconciliação, além de persistência local e central. Operar sem rede depende também de relógio confiável, agenda válida, calibração e prontidão física. \cite{DC_COMPILADO,DC_ANALISE}
\section{Relógio de tempo real e alimentação}
Um RTC com alimentação de reserva pode manter a contagem de tempo quando a alimentação principal é interrompida. Isso não alimenta a Raspberry Pi ou os motores. As variantes DS3231 e DS3231M não devem ter suas especificações confundidas; a documentação fotográfica do projeto menciona DS3231M em placa HW-084. O circuito do módulo e o tipo de bateria precisam ser identificados antes da instalação. \cite{ADI_DS3231M,DC_COMPILADO}
Capacidade nominal de uma fonte não corresponde ao consumo do equipamento. Corrente, queda de tensão, aquecimento e comportamento nos transitórios devem ser medidos na montagem final; a previsão de 5 V/10 A não demonstra suficiência elétrica por si só. \cite{DC_CIRCUITOS}
\section{Aplicativo, servidor e autenticação}
A separação entre interface, API e execução local permite identificar quem valida uma ação e quem a realiza fisicamente. No DoseCerta, Care utiliza React/TypeScript; a API FastAPI intermedeia as regras de cuidado e a coordenação com o dispositivo. O material descreve JWT e Supabase, mas a configuração efetiva deve ser vinculada à versão demonstrada. \cite{DC_APP}
\section{Acessibilidade e interação}
As WCAG 2.2 oferecem critérios para conteúdo web perceptível, operável, compreensível e robusto. Entre os pontos pertinentes estão contraste, uso por teclado, identificação de controles e não depender exclusivamente de cor. Referenciar as diretrizes não equivale a atestar conformidade; é preciso delimitar páginas, estados e critérios efetivamente avaliados. \cite{WCAG22}
\section{Visão computacional e interpretação de eventos}
A arquitetura reserva a visão para a observação associada à tomada. Uma pontuação produzida por algoritmo não deve ser interpretada automaticamente como probabilidade clínica calibrada. O método, o conjunto de dados, os limiares e as métricas específicos do DoseCerta serão descritos pelo responsável no capítulo 12. \cite{DC_COMPILADO}
\section{Agentes de IA e ferramentas}
O agente conversacional interpreta solicitações e pode consultar contexto ou solicitar ações por ferramentas. Instruções ao modelo, regras do servidor e mecanismos de aprovação são camadas distintas. No projeto, a documentação do Foundry separa o agente de voz e a cópia com MCP, pois o fluxo de voz ainda não atende aos pedidos de aprovação de escrita. \cite{DC_FOUNDRY}
\section{Prototipagem e limitações de materiais}
A aquisição de PLA não comprova o material de cada peça, nem a adequação sanitária de superfícies em contato com medicamentos. O processo de impressão, o acabamento, a limpeza, o desgaste e a estabilidade devem ser documentados para a configuração ensaiada. Não se declara biocompatibilidade, esterilizabilidade ou certificação do protótipo com base apenas na denominação do polímero. \cite{DC_COMPILADO,DC_LISTA}
